\section{Halo integrals}
\label{sec: halo integrals}

The astrophysical dependence of the event rate in direct detection experiments (see section~\ref{sec: direct detection}) comes from the DM velocity distribution and density in the Solar neighbourhood. For the case of standard interactions, the \emph{halo integral} encodes the local DM velocity distribution dependence of the event rate  and is defined as
\begin{equation}
    \eta (v_{\rm min}, t) \equiv \int_{v>v_{\rm min}} d^{3}v\, \frac{\tilde{f}_{\rm det}(\mathbf{v},t)}{v} \; ,
    \label{eq: halo integral}
\end{equation}
%
where $\mathbf{v}$ is the relative velocity between the DM and the target nucleus or electron in the detector, with $v=|\mathbf{v}|$, $\tilde{f}_{\rm det}({\bf v},t)$ is the local DM velocity  distribution in the detector reference frame, and  $v_{\rm min}$ is the minimum  speed required for the DM particle to impart a recoil energy and momentum in the detector (given in eqs.~\eqref{eq: vmin} and \eqref{eq: vmin-e} for nuclear and electron recoils, respectively). Determining the influence of the LMC on the halo integrals in the Solar region directly reflects the expected change in direct detection event rates.

 

We extract the halo integrals of the MW-LMC analogues by boosting the local DM velocity distribution of each halo from the Galactic reference frame to the detector frame,
\begin{equation}
\tilde{f}_{\rm det}({\bf v},t) = \tilde{f}_{\rm gal}({\bf v}+{\bf v}_s+{\bf v}_e(t))\; , 
\end{equation}
where ${\bf v}_e(t)$ is the Earth's velocity with respect to the Sun, ${\bf v}_s={\bf v}_c +{\bf v}_{\rm pec}$ is the Sun's velocity in the Galactic rest frame, ${\bf v}_c$ is the Sun's circular velocity, and ${\bf v}_{\rm pec}=(11.10, 12.24, 7.25)$~km/s \cite{Schoenrich:2009bx} is the peculiar velocity of the Sun in Galactic coordinates with respect to the Local Standard of Rest. To boost the DM velocity distribution to the detector rest frame, we take $|{\bf v}_c|=v_c=220$~km/s. For simplicity, we neglect the small eccentricity of the Earth's orbit. In the following, we present the time-averaged halo integrals, which are averaged over one year.



In figure~\ref{fig: halo panel lmc candidate} we present the time-averaged halo integrals as a function of $v_{\rm min}$  in the Solar region for the best fit Sun's position for the same four halos whose local DM speed distributions are shown in figure~\ref{fig: fv lmc candidates}: halos 2 (top left), 6 (top right), 13 (bottom left) and 15 (bottom right), for the
snapshot closest to the LMC's pericenter approach. The black and red solid lines are the halo integrals computed from the mean value of the velocity distributions of the DM particles originating from the MW+LMC and the MW only, respectively. The shaded bands correspond to the $1\sigma$ uncertainties in the halo integrals and are obtained from the DM velocity distribution at one standard deviation from the mean. The panels below the halo integral plots show the relative difference between the MW+LMC and the MW-only halo integrals, defined as $(\eta_{\rm MW+LMC}-\eta_{\rm MW})/\eta_{\rm MW}$.


\begin{figure}
\centering
\includegraphics[width=0.8\textwidth]{graphics/halo_panel_ratio.pdf}
 \caption{Time-averaged halo integrals for halos 2 (top left), 6 (top right), 13 (bottom left) and 15 (bottom right) in the Solar region for the best fit Sun's position, for the snapshot closest to the LMC's pericenter approach. The black and red curves show the halo integrals for the DM particles originating from the MW+LMC and MW only, respectively. In each case, the solid lines and the shaded bands correspond to the halo integrals obtained from the mean DM velocity distribution and the DM velocity distribution at $1\sigma$ from the mean, respectively. The value of $\kappa_{\rm LMC}$ is also specified on each panel. The panels below the halo integral plots show the relative difference between the MW+LMC and the MW-only halo integrals, $(\eta_{\rm MW+LMC}-\eta_{\rm MW})/\eta_{\rm MW}$.}
 \label{fig: halo panel lmc candidate}
\end{figure}

As seen in figure~\ref{fig: halo panel lmc candidate}, halos 6 and 15 show some differences in the tails of the halo integrals of the MW+LMC and the MW-only, with their relative difference reaching $\sim 6$ for halo 6 and  $\sim 0.5$ for halo 15. The halo integrals of halos 2 and 13 do not show any visible deviations between the MW+LMC and the MW-only, and their relative differences are  smaller than 0.1 for halo 2 and 0.01 for halo 13. This is despite the fact that halo 13 has a higher $\kappa_{\rm LMC}$ in the Solar region compared to the other three halos. This highlights  the importance of the particular shape and peak speed of the LMC's speed distribution in the detector reference frame, in the Solar region of each MW analogue.






To quantify the changes in the tails of the halo integrals of the native DM particles of the MW and the total DM particles originating from MW+LMC, we define a dimensionless metric, 
\begin{equation}
\Delta \eta =\sum_{v_{\rm min}^i\geq 0.7 v_{\rm esc}^{\rm det}}\left[\eta_{\rm MW+LMC}(v_{\rm min}^i)-\eta_{\rm MW}(v_{\rm min}^i)\right]\Delta v_{\rm min}\; ,
    \label{eq: delta eta}
\end{equation}
where $\Delta v_{\rm min}$ is the bin size in $v_{\rm min}$, and $v_{\rm min}^i$ denotes the midpoint of the bins in $v_{\rm min}$ at which the halo integrals of the MW+LMC, $\eta_{\rm MW+LMC}$, and MW only, $\eta_{\rm MW}$, are evaluated. The sum runs over all bins with $v_{\rm min}^i$ larger than 70\% of the local escape speed from the MW in the detector rest frame, $v_{\rm esc}^{\rm det}$, which is estimated from the largest $v_{\rm min}$ where $\eta_{\rm MW}$ is nonzero. The values of $v_{\rm esc}^{\rm det}$ in the Solar region for the best fit Sun's position for the 15 MW-LMC analogues are given in table~\ref{tab:Localrho}, for the simulation snapshot closest to the LMC’s pericenter approach.



The metric in eq.~\eqref{eq: delta eta} reflects the changes in the exclusion limits set by direct detection experiments for the MW+LMC and MW-only distributions at low DM masses. Since the integral in eq.~\eqref{eq: halo integral} is computed for speeds greater than $v_{\rm min}$, and $v_{\rm min}$ depends inversely on the DM mass (eqs.~\eqref{eq: vmin} and \eqref{eq: vmin-e}), the exclusion limits in direct detection experiments become sensitive to small changes in the high speed tail of the halo integrals for low DM masses. Consequently, $\Delta\eta$ was defined to include only the differences in the halo integrals for $v_{\rm min}$ larger than 70\% of $v_{\rm esc}^{\rm det}$ to numerically reflect the variations in the tail of the halo integral and direct detection exclusion limits  at low DM masses. 
We have checked various other metrics for $\Delta\eta$, including the relative difference, the difference in the area under the curves, and considering different fractions of $v_{\rm esc}^{\rm det}$ in the above metric, with all showing similar general trends. The current definition 
preserves the global trends while providing the most intuitive connection between the halo integral plots and the direct detection exclusion limits (presented in section~\ref{sec: direct detection}) calculated therefrom.  


In our analysis we find three key factors that contribute to changes in the tail of the halo integrals: 1) the percentage of DM particles originating from the LMC in the Solar region, 2) the Sun's position (and hence the \emph{Solar region}) in the simulations, and 3) the MW response due to the motion of the LMC as it traverses its orbit near pericenter. In the following sections we discuss how the results depend on each of these phenomena in detail.


%%%%%%%%%%%%%%%%%%%%%%%%%%%%%%%%%%%%%%%%%%%%%%%%%%%%%
\subsection{Impact of the DM particles originating from the LMC}
\label{sec: percent LMC}

\begin{figure}
     \centering
     \begin{subfigure}[b]{0.49\textwidth}
         \includegraphics[width=\textwidth]{graphics/delta_eta_v_percent_lmc_candidates.png}
     \end{subfigure}
    \hspace{-0.25cm}
     \begin{subfigure}[b]{0.49
     \textwidth}
         \includegraphics[width=\textwidth]{graphics/70eta_percent_v_delta_snapshots.png}
     \end{subfigure}
        \caption{The correlation between the change in the tail of the halo integrals due to the LMC particles, $\Delta\eta$, and $\kappa_{\rm LMC}$ for three Solar regions: the Solar region for the best fit Sun's position (black squares), the Solar region that maximizes $\Delta\eta$ (yellow dots), and the Solar region that minimizes the $\Delta\eta$ (blue dots). Left panel: all 15 MW-LMC analogues at the snapshot closest to the LMC's pericenter approach. Right panel: different snapshots of halo 13, ranging from $\sim 314$~Myr before the present day to $\sim 175$~Myr after.}
        \label{fig: eta v percent}
\end{figure}



The DM particles originating from the LMC with high enough speeds become unbound to the LMC at infall, with some number at any given time ending up in the Solar region of the MW. Since these particles have on average higher speeds than the native DM particles of the  MW, it is expected that they will affect the high speed tail of the halo integrals, with a higher value of $\kappa_{\rm LMC}$ contributing to a more pronounced effect. 

Figure~\ref{fig: eta v percent} shows the correlation between the percentage of DM particles originating from the LMC in the Solar region, $\kappa_{\rm LMC}$, and the change in the tail of the halo integral due to the LMC particles, $\Delta\eta$ (eq.~\eqref{eq: delta eta}), for the best fit Sun's position (black squares), the Solar region that maximizes $\Delta\eta$ (yellow dots), and the Solar region that minimizes $\Delta\eta$ (blue dots). The left panel shows the results for the 15 MW-LMC analogues at the snapshot closest to the pericenter approach of the LMC, while the right panel is for different snapshots of halo 13, ranging from $\sim 314$~Myr before the present day to $\sim 175$~Myr after. In both panels, the $\Delta\eta$ for the best fit Sun's position is in general close to the maximum $\Delta \eta$, both showing an increase with $\kappa_{\rm LMC}$ in the Solar region. The minimum $\Delta \eta$ is zero or close to zero for a number of MW-LMC analogues and for some snapshots of halo 13, but still shows an increase with $\kappa_{\rm LMC}$ for halo 13. As discussed in section~\ref{sec: DM density}, $\kappa_{\rm LMC}$ generally increases with $M_{\rm Infall}^{\rm LMC}/M_{200}^{\rm MW}$. Therefore, a higher LMC to MW mass ratio would in general result in a larger $\Delta \eta$, depending on the particular Sun's position considered.




To better visualize the variation of $\Delta \eta$ in halo 13, in figure~\ref{fig: bestfit deltaeta v snapshot} we  present
$\Delta\eta$ in the Solar region for the best fit Sun's positions for different snapshots as a function of the snapshot time relative to the present day snapshot, $t-t_{\rm Pres.}$. The colour bar specifies the range of $\kappa_{\rm LMC}$. As expected,  snapshots with the highest value of $\kappa_{\rm LMC}$ near LMC's pericenter approach\footnote{The two snapshots occurring at $\sim 174$ and $\sim 191$~Myr before the present day snapshot are tied for the highest value of $\kappa_{\rm LMC}$ with 1.0\% each. From those, the one closest to the pericenter snapshot ($\sim 41$~Myr before LMC's pericenter approach) has the highest $\Delta\eta$.} also have the highest $\Delta\eta$, with snapshots far from pericenter dropping off in  both $\kappa_{\rm LMC}$ and $\Delta\eta$. 



\begin{figure}[t]
    \centering
    \includegraphics[scale=.7]{graphics/delta_eta_v_delta_t.pdf}
    \caption{$\Delta \eta$ in the Solar region for the best fit Sun's positions for different snapshots in halo 13, plotted against the snapshot time relative to the present day snapshot, $t-t_{\rm pres}$. The snapshot times range from $\sim 314$~Myr before the present day to $\sim 175$~Myr after. The colour bar shows the range of $\kappa_{\rm LMC}$. The snapshots for the LMC's pericenter approach (Peri.), present day (Pres.) and $\sim 175$~Myr after the present day (Fut.) are specified with vertical dashed black lines.}
    \label{fig: bestfit deltaeta v snapshot}
\end{figure}



Notwithstanding the relationship between $\Delta\eta$ and $\kappa_{\rm LMC}$, there remains a scatter in the values of $\Delta\eta$ for systems with equal or similar values of $\kappa_{\rm LMC}$,  due to the particular choice of Sun's position for specifying the Solar region. This can be seen in both panels of figure~\ref{fig: eta v percent}, where there are large differences between the minimum and maximum $\Delta \eta$ for  the same or similar values of $\kappa_{\rm LMC}$. 
This leads us to consider not just the impact of $\kappa_{\rm LMC}$ on $\Delta \eta$, but also the effect of the exact Sun-LMC geometry, and whether the best fit Sun's position is a privileged position with respect to maximizing $\Delta\eta$. We explore this in the next section. 





\subsection{Variation due to the Sun-LMC geometry}
\label{sec: Variation due to Sun-LMC Geometry}


For the 15 MW-LMC analogues we find a degree of variation in the values of $\kappa_{\rm LMC}$  due to the choice of the Solar region. In particular, $\kappa_{\rm LMC}$ can vary at most by a factor of $\sim 2$ depending on the MW-LMC analogue (see e.g.~the last column of table~\ref{tab:Localrho}). However, as discussed in section~\ref{sec: percent LMC}, for Solar regions with similar values of $\kappa_{\rm LMC}$ there is a large scatter in how much the tails of the MW+LMC halo integrals can deviate from their MW only counterparts.  This can also be seen in figure~\ref{fig: halo panel lmc candidate}, where the largest deviation in the tail of the halo integral is seen for halo 6 with $\kappa_{\rm LMC}=0.038\%$, while halo 13, with the highest $\kappa_{\rm LMC}$ of 2.3\%, shows a very small variation. 
This implies that the  value of $\kappa_{\rm LMC}$ is not the only important factor in specifying $\Delta \eta$, but  the particular Sun-LMC geometry of the chosen Solar region is similarly important.%



Figure~\ref{fig: halo panel bestfit} shows the time-averaged halo integrals for the MW+LMC (black) and the MW only (red) DM populations for the present day snapshot of halo 13 for two different Solar regions: the best fit Sun's position (left panel) and the Solar region that minimizes $\Delta\eta$ (right panel). The value of $\kappa_{\rm LMC}$ and the cosine angles corresponding to the particular Sun's position (eq.~\eqref{eq: cosine angles}) are also specified in each panel. The panels below the halo integral plots show the relative difference between the MW+LMC and the MW-only halo integrals. Clear differences between the tails of the MW+LMC and the MW halo integrals are visible in the case of the best fit Sun's position, while the variations in the halo integral tails are small for the Solar region that minimizes $\Delta \eta$. This figure highlights the importance of the Sun-LMC geometry, and that the same snapshot with similar values of $\kappa_{\rm LMC}$  can differ greatly in the tails of the halo integrals due to the choice of the Solar region. 


To differentiate the effects on $\Delta \eta$ due to the choice of the Sun-LMC geometry from the effect of $\kappa_{\rm LMC}$, we study the cosine angles  that parameterize the Sun-LMC geometry, as given in eq.~\eqref{eq: cosine angles}. Figure~\ref{fig: cosine angles and sun position} shows the allowed Sun's positions in the parameter space of the two cosine angles for halo 13 at the present day snapshot. The colour bars show the range of $\kappa_{\rm LMC}$ and $\Delta\eta$ in the left and right panels, respectively. The black square in each panel shows the observed values of the cosine angles from eq.~\eqref{eq: best fit cosine angles}. Comparing the two panels of the figure shows that $\Delta\eta$ maximizes in the quadrant where $\cos \alpha$ and $\cos \beta$ are negative, despite $\kappa_{\rm LMC}$ reaching its maximum in the positive $\cos \alpha$ and $\cos \beta$ quadrant. 

\begin{figure}
     \centering
         \includegraphics[width=0.85\textwidth]{graphics/halo13_panel_ratio.pdf}
        \caption{Time-averaged halo integrals for halo 13 at the present day snapshot for the MW+LMC (black) and the MW only (red) DM populations for the  best fit Sun's position (left panel) and the Solar region that has the minimum $\Delta\eta$ (right panel). The solid lines and the shaded bands correspond to the halo integrals obtained from the mean DM velocity distribution and the DM velocity distribution at $1\sigma$ from the mean, respectively. The value of $\kappa_{\rm LMC}$ and the cosine angles (eq.~\eqref{eq: cosine angles}) are also specified on each panel. The panels below the halo integral plots show the relative difference between the MW+LMC and the MW-only halo integrals.}
        \label{fig: halo panel bestfit}
\end{figure}



The value of $\kappa_{\rm LMC}$ varies on average by $0.15\%$ between different allowed Sun's positions of a given snapshot, while it can vary by up to a percent between different snapshots. Hence, within a snapshot the dominant factor that impacts $\Delta\eta$ is the particular Sun-LMC geometry. Furthermore, we find that across snapshots $\Delta\eta$ tends to have its maximum values in the quadrants where  $\cos \beta$ is negative. A negative $\cos \beta$ indicates that the velocity vector of the LMC analogue is in the opposite direction of the Sun's velocity vector, leading to larger relative speeds of the particles originating from the LMC with respect to the Sun, and thus resulting in a larger $\Delta \eta$. Since  the observed cosine angles are also negative, this implies that variations in the tail of the halo integral for the best fit Sun's position should be close to the maximum possible variation from other allowed Sun's positions. This can also be seen from figure~\ref{fig: eta v percent}, where the values of $\Delta \eta$ for the best fit Sun's positions (black squares) are close to the maximum $\Delta \eta$ (yellow dots). We can therefore conclude that the best fit Sun's position is indeed in a privileged position with respect to maximizing $\Delta\eta$, by virtue of the observed $\cos \beta$ being negative, i.e.~the Sun's velocity and the LMC's velocity being predominantly in opposite directions. As a consequence, for the actual MW we expect the LMC to maximally affect the tail of the halo integral.


\begin{figure}[t]
    \centering
    \includegraphics[width=\textwidth]{graphics/153_cosine_panels_70vmin.pdf}
    \caption{
    Cosine angles that parameterize the Sun-LMC geometry  (eq.~\eqref{eq: cosine angles})  for halo 13 at the present day snapshot for all allowed Sun's positions coloured by the value of $\kappa_{\rm LMC}$ (left panel) and $\Delta\eta$ (right panel). The observed values of the cosine angles (eq.~\eqref{eq: best fit cosine angles}) is specified with a black square on each panel.
  }
    \label{fig: cosine angles and sun position}
\end{figure}

We have also checked the dependence of $\Delta \eta$  on the size of our defined Solar region for the present day snapshot of halo 13. We find that decreasing the size of the Solar region leads to a significant increase in $\Delta \eta$, due to better sensitivity to the best fit Sun's position. In particular, decreasing the opening angle of the cone from $\pi/4$ to $\pi/6$ while keeping the spherical shell width the same, increases $\Delta \eta$ by $18\%$. Decreasing the shell width to $7-9$~kpc while keeping the opening angle of the cone the same has a less significant effect and increases $\Delta \eta$ by $7\%$. Decreasing both the opening angle of the cone to $\pi/6$ and the shell width to $7-9$~kpc, increases $\Delta \eta$ by $78\%$. This comes at the cost of significantly reducing the number of DM particles from the MW and LMC in the Solar region (as discussed in section~\ref{sec: DM density}), and leading to very large Poisson uncertainties. Our results are therefore conservative with respect to the choice of the Solar region. Increasing the size of the Solar region, on the other hand, results in only a slight decrease in $\Delta \eta$, due to losing sensitivity to the best fit Sun's position. In particular, increasing the opening angle of the cone from  $\pi/4$ to $\pi/2$, while keeping the spherical shell width the same, decreases $\Delta \eta$ by only $2\%$.  

\subsection{MW response to the LMC}
\label{sec: MW response}

In addition to the particles originating from the LMC in the Solar region, the response of the local DM halo of the MW to the LMC's orbit can cause variations in the high speed tail of the local DM velocity distribution. The MW response to the LMC has been observed and studied before in idealized simulations~\cite{Besla:2019xbx, Donaldson:2021byu}, but it is important to test it in a fully cosmological setting where halos have multiple accretion events over their formation history. In this section, we explore the effect of this response on the tail of the halo integral in our cosmological simulations, and distinguish it from the effect of the high speed DM particles in the Solar region that originate from the LMC.



Figure~\ref{fig: halo int response} shows the time-averaged halo integrals for the four representative snapshots in halo 13: the isolated MW analogue (Iso.), the LMC's pericenter approach (Peri.), the present day MW-LMC analogue (Pres.), and the future MW-LMC analogue (Fut.). The halo integrals of the three latter snapshots are shown for the MW+LMC (solid coloured curves) and the MW only (dashed coloured curves) for the best fit Sun's position. The isolated MW snapshot has no LMC-like satellite and hence its halo integral (solid black curve) is extracted from the DM particles of the MW in a spherical shell with radii between 6 and 10 kpc from the Galactic center. For comparison, the blue curve shows the halo integral obtained from a Maxwellian velocity distribution with a peak speed of 220~km/s and truncated at the escape speed of 544~km/s from the Galaxy, as is commonly assumed in the SHM.


\begin{figure}[t]
    \centering
    \includegraphics[scale=0.55]{graphics/rescaled_MW_response_plot.pdf}
    \caption{Time-averaged halo integrals for four  snapshots in halo 13: the isolated MW analogue (Iso.), the LMC's pericenter approach (Peri.), the present day MW-LMC analogue (Pres.), and the future MW-LMC analogue (Fut.). For each snapshot, the solid/dashed  lines and the shaded bands correspond to the halo integrals obtained from the mean DM velocity distribution and the DM velocity distribution at $1\sigma$ from the mean, respectively. For the present day, pericenter, and future snapshots the halo integrals are presented in the best fit Solar region
    for the MW+LMC (solid coloured curves) and MW-only (dashed coloured curves) DM populations. The isolated MW snapshot has no LMC analogue, so its MW halo integral (solid black curve) is shown for a Solar region defined as a spherical shell with radii between 6 and 10 kpc from the Galactic center. The solid blue curve shows the SHM halo integral obtained from a Maxwellian velocity distribution with a peak speed of 220~km/s and truncated at the escape speed of 544~km/s from the Galaxy.
   }
    \label{fig: halo int response}
\end{figure}

 

A comparison of the halo integral of the isolated MW  with  the MW-only halo integrals at the three other snapshots  shows how the native DM particles of the MW in the Solar region are boosted in response to the passage of the LMC. The isolated MW snapshot occurs $\sim 2.8$~Gyr before the present day snapshot and is a proxy for the  MW in the absence of the LMC's influence. The halo integral for this snapshot (solid black curve) is closest to the SHM halo integral, although its tail is slightly more extended, reaching $v_{\rm min} \sim 800~\rm{km/s}$. As the LMC reaches its first pericenter approach, the tail of the halo integral for the native DM population of the MW (dashed green curve) is boosted to $v_{\rm min} \sim 900~\rm{km/s}$. Since the present day LMC is not too far from its pericenter approach, the tail of the halo integral at the present day snapshot (dashed orange curve) shows a comparable boost to the pericenter snapshot. 
Finally, the tail of the halo integral for the local DM population of the MW returns to $v_{\rm min}  \sim 800$~km/s at the future MW-LMC snapshot (dashed magenta curve), which occurs $\sim 175$~Myr after the present day snapshot\footnote{This period of time is comparable to the overdensity wake induced by the passage of a satellite galaxy in a host DM halo in the Auriga simulations~\cite{Gomez2016}.}. 



The addition of DM particles originating from the LMC in the Solar region further shifts the tails of the halo integrals to higher speeds.
In particular, the pericenter snapshot has the highest  $\kappa_{\rm LMC}$ of 0.85\% and also shows the highest boost in the tail of its MW+LMC halo integral (solid green curve), reaching $v_{\rm min} \sim 1000~\rm{km/s}$, which is $\sim100~\rm{km/s}$ higher than the reach of its MW-only counterpart. Similarly, with $\kappa_{\rm LMC}=0.26$\%, the present day snapshot has a  MW+LMC halo integral which exhibits a large difference compared its MW-only counterpart in the high speed tail. The future MW-LMC snapshot 
has $\kappa_{\rm LMC}=0.22$\%, and the tail of its  MW+LMC halo integral is boosted by $\sim50~\rm{km/s}$ compared to its MW-only counterpart. 

Comparing the boost of the native DM population of the MW in the pericenter and present day snapshots to the boost due to the presence of DM particles originating from the LMC reveals that the impact on the tail of the halo integral is of the same order of magnitude. 
Figure~\ref{fig: halo int response} demonstrates that the DM particles in the Solar neighborhood can be boosted from  $v_{\rm min} \sim 800~\rm{km/s}$ in the absence of the LMC (solid black curve) to more than $v_{\rm min} \sim 950~\rm{km/s}$ at the present day (solid orange curve), a combined increase of greater than $\sim 150~\rm{km/s}$ due to the MW response and the presence of high speed LMC particles in the Solar region. 





